% !Mode:: "TeX:UTF-8"
\chapter{实验与分析}

\section{实验目的}
为验证本文提出的面向博物馆导览的对话式数字人展示设计是否能够回应前文提出的研究问题，本文开展用户实验，对对话式数字人导览原型系统在实际导览任务中的表现进行分析。实验重点围绕展品知识讲述支持、导览情境响应以及注意力引导三个方面展开。

具体而言，本文希望通过实验回答以下三个问题：

\textbf{（1）}系统能否在导览过程中有效承载展品知识，并支持围绕展品内容展开持续、连贯的讲述。该部分主要对应本文在 Conversation 维度中提出的分层讲解、追问延续与知识组织设计。

\textbf{（2）}系统能否在连续导览过程中维持并利用当前导览情境，根据用户当前关注的展品对象与前序交流内容作出相对合理的响应。该部分主要对应本文在 Context 维度中提出的情境状态维护与导览过程适配机制。

\textbf{（3）}数字人能否通过语音、动作、视线等表达方式形成较明确的注意力引导，使用户能够更容易判断当前应关注的展品对象与讲解重点。该部分主要对应本文在 Embodied Guidance 维度中提出的表达映射与多模态协同设计。

基于上述目标，本文通过对比不同导览方式下的用户体验与过程表现，考察对话式数字人导览原型系统在知识讲述、情境响应与注意力引导方面的实际效果，从而为本文提出的对话式数字人展示设计提供验证依据。

\section{实验设计}
为了评估本文所提出的对话式数字人展示设计在实际导览任务中的表现，本研究采用对比实验方法，将本文实现的 CCEG 对话式数字人导览原型系统与单轮问答式数字人导览基线系统进行比较，从而观察不同导览机制下用户体验与交互行为的差异。本文的研究目标在于验证 CCEG 设计作为整体导览方案相较于常见问答式数字人导览方式的有效性，本研究仅设置一个基线条件，不进一步拆分各模块进行消融比较。

实验设置两个条件：

\textbf{（1）基线条件（单轮问答式数字人导览系统）}
基线条件下，系统使用与实验系统相同的展品知识来源、数字人形象与语音输出方式，但仅针对用户当前问题给出单次回答，不维护连续对话中的导览情境状态，不进行主动追问、主题推进与讲解层次调节，也不提供面向展品对象的专门动作、视线与语气引导策略。该条件用于模拟当前较为常见的问答式数字人导览形态。

\textbf{（2）实验条件（CCEG系统）}
实验条件采用本文设计实现的对话式数字人导览原型系统。该系统在相同知识基础与数字人载体上，进一步加入面向导览任务的对话组织、情境状态维护与具身引导机制，使系统能够围绕当前展品持续组织讲解，并结合参观进程与表达策略进行导览引导。

实验选取同一展区中的若干典型展品作为导览对象。为保证比较的可比性，两种条件下使用相同的展品内容与任务要求。参与者在实验过程中需要围绕这些展品完成一次简短的导览体验任务，包括浏览展品信息、向系统提出问题以及根据系统反馈继续了解相关内容。

本实验共邀请 12 名参与者，参与者均为具有基础博物馆参观经验的学生或普通观众。为减弱体验顺序可能带来的影响，实验采用 A/B 两组交叉验证的方式进行：将参与者平均分为 A、B 两组，A 组先体验基线条件、后体验实验条件，B 组先体验实验条件、后体验基线条件。参与者在每次体验结束后分别填写问卷，对导览过程中的讲解连贯性、情境响应以及注意力引导等方面进行评价。

实验流程如下：
\begin{enumerate}
    \item 向参与者介绍实验目的与基本操作方式；
    \item A 组先体验基线条件、B 组先体验实验条件，并分别完成对应导览任务；
    \item 参与者填写第一轮体验问卷并进行简短反馈；
    \item A 组继续体验实验条件、B 组继续体验基线条件，并重复上述流程。
\end{enumerate}

通过对比两种条件下的用户评价与过程表现，可以观察完整 CCEG 导览设计相较于单轮问答式数字人导览方式，是否在知识讲述支持、情境响应能力以及注意力引导方面表现更优。

\section{评估指标设计}
为了对应验证本文提出的三个研究问题，实验评估不再采用笼统的整体满意度划分，而是围绕知识讲述支持、情境响应能力和注意力引导效果三个维度展开。该指标设计分别对应前文中提出的展品知识交互化呈现问题、导览情境维持问题以及数字人注意力引导不足问题，并与 CCEG 框架中的 Conversation、Context 和 Embodied Guidance 三个部分保持一致。

本研究采用问卷评价的方式进行分析。其中，问卷部分采用 5 点 Likert 量表，1 表示“非常不同意”，5 表示“非常同意”。

\textbf{（1）知识讲述支持}

该维度用于评估系统是否能够在导览过程中有效承载展品知识，并支持围绕同一展品展开持续、连贯的讲述。前文在对话组织设计中已提出，系统应通过分层叙事单元、追问路径与回收策略，使导览对话能够围绕观众理解进程逐步展开，而不是停留在碎片化问答层面。

基于此，本文设置如下评价条目：
\begin{itemize}
    \item 系统能够围绕同一展品持续展开讲解。
    \item 系统的回答能够与我的提问保持连贯，而不是零散的信息。
    \item 系统提供的信息有助于我形成对展品内容的整体理解。
\end{itemize}

这一维度主要用于观察系统在知识组织、讲述延续和内容连贯性方面的表现。

\textbf{（2）情境响应能力}

该维度用于评估系统能否在连续导览过程中维持并利用当前导览情境，根据用户当前关注的展品对象和前序交流内容作出相对合理的响应。前文在情境状态设计中指出，系统需要围绕导览对象、参观进程和交流历史维护情境状态，以支持内容选择和策略调节。

基于此，本文设置如下评价条目：
\begin{itemize}
    \item 系统的回答能够贴合我当前关注的展品。
    \item 在连续对话中，系统能够记住前面的交流内容。
    \item 系统的讲解能够随着参观进程的变化进行调整。
\end{itemize}

这一维度主要用于考察系统在情境维持、上下文利用和导览过程适配方面的能力。

\textbf{（3）注意力引导效果}

该维度用于评估数字人能否通过语音、动作、视线等表达方式形成相对明确的注意力引导，使用户能够判断当前应关注的展品对象与讲解重点。前文已指出，具身数字人的关键并不只是“能说”，而是要通过讲解语义与行为表达的联动，使非语言表达与空间对象、讲解意图保持一致，从而形成可感知的注意力引导。

基于此，本文设置如下评价条目：
\begin{itemize}
    \item 我能够感受到数字人在主动引导我关注展品。
    \item 数字人的动作、视线或语气有助于我判断当前讲解重点。
    \item 数字人的表达更像导览引导，而不仅是信息播报。
\end{itemize}

这一维度主要用于观察系统在注意力组织、重点提示和引导可感知性方面的表现。


\section{实验结果分析}
本节采用定量统计方法对实验结果进行分析。分析数据主要包括两类：一是两种导览条件下各问卷条目的评分结果，二是按知识讲述支持、情境响应能力和注意力引导效果三个维度汇总后的总体得分。对上述数据分别统计均值与标准差，用于描述不同导览条件下各评价维度的总体表现。

由于本实验采用同一参与者分别体验两种条件的 A/B 两组交叉验证设计，后续结果分析以配对样本数据为基础展开。若样本分布近似满足正态性假设，则采用配对样本 t 检验比较两种条件下各维度评分的差异；若数据不满足正态性要求，则采用 Wilcoxon 符号秩检验作为非参数替代方法。统计检验采用 95\% 置信水平，显著性水平设定为 $p < 0.05$。
